\usemodule[memos][fontsize=16pt,fontstyle=ss]
\usemodule[basicexam]
\usemodule[poriginal][color=red]
\starttext

\TitlePage[title=这是文章标题,subtitle=这是副标题,author=\moduleparameter{poriginal}{color},]{}

\Topics{这是篇章目录页}

\Topic{这是主要篇章}

  \Subject{这是主题}  \currentmoduleparameter{color} \input knuth

\Topic{Font}

  \Subject{FontFamily}
  
\startxtable
  \startxrow
    \startxcell fontfamily     \stopxcell
    \startxcell groupedcommand \stopxcell
    \startxcell declaredcommand\stopxcell
    \startxcell result         \stopxcell
  \stopxrow
  \startxrow
    \startxcell Serif          \stopxcell
    \startxcell \type{\style[rm]{Abc}}\stopxcell
    \startxcell \type{\rm Abc} \stopxcell
    \startxcell {\rm Abc}      \stopxcell
  \stopxrow
  \startxrow
    \startxcell Sans           \stopxcell
    \startxcell \type{\style[ss]{Abc}}\stopxcell
    \startxcell \type{\ss Abc} \stopxcell
    \startxcell {\ss Abc}      \stopxcell
  \stopxrow
  \startxrow
    \startxcell Mono           \stopxcell
    \startxcell \type{\style[tt]{Abc}}\stopxcell
    \startxcell \type{\tt Abc} \stopxcell
    \startxcell {\tt Abc}      \stopxcell
  \stopxrow
  \startxrow
    \startxcell Hand           \stopxcell
    \startxcell \type{\style[hw]{Abc}}\stopxcell
    \startxcell \type{\hw Abc} \stopxcell
    \startxcell {\hw Abc}      \stopxcell
  \stopxrow
\stopxtable

\startxtable
  \startxrow
    \startxcell fontstyle      \stopxcell
    \startxcell groupedcommand \stopxcell
    \startxcell declaredcommand\stopxcell
    \startxcell result         \stopxcell
  \stopxrow
  \startxrow
    \startxcell normal          \stopxcell
    \startxcell \type{\style[tf]{Abc}}\stopxcell
    \startxcell \type{\tf Abc} \stopxcell
    \startxcell {\tf Abc}      \stopxcell
  \stopxrow
  \startxrow
    \startxcell Bold           \stopxcell
    \startxcell \type{\style[bold]{Abc}}\stopxcell
    \startxcell \type{\bf Abc} \stopxcell
    \startxcell {\bf Abc}      \stopxcell
  \stopxrow
  \startxrow
    \startxcell italic          \stopxcell
    \startxcell \type{\style[it]{Abc}}\stopxcell
    \startxcell \type{\it Abc} \stopxcell
    \startxcell {\it Abc}      \stopxcell
  \stopxrow
  \startxrow
    \startxcell CAPS           \stopxcell
    \startxcell \type{\style[sc]{Abc}}\stopxcell
    \startxcell \type{\sc Abc} \stopxcell
    \startxcell {\sc Abc}      \stopxcell
  \stopxrow
\stopxtable
  
  \Subject{Font Decoration}
  
\useURL[bing]   [http://bing.com/]          [] [I prefer dog.]
\useurl[mail]   [mailto:nobody@example.zzz] [] [visible@mailaddress.zzz]

\startpoints
  \item \hbox to 8em{GOTO LINK \hfil} \goto{Wiki}[url(http://wiki.contextgarden.net)]
  \item \hbox to 8em{GOTO PAGE \hfil} \goto{Other page}[page(3)]
  \item \hbox to 8em{URL       \hfil} \from[bing]
  \item \hbox to 8em{MAIL      \hfil} \from[mail]
\stoppoints
  
\startpoints[horizontal,two]\ssx
  \startitem \type{\overbar{Lorem ipsum}       }\hfill\overbar{Lorem ipsum}       \stopitem
  \startitem \type{\underbar{Lorem ipsum}      }\hfill\underbar{Lorem ipsum}      \stopitem
  \startitem \type{\overstrike{Lorem ipsum}    }\hfill\overstrike{Lorem ipsum}    \stopitem
  \startitem \type{\understrike{Lorem ipsum}   }\hfill\understrike{Lorem ipsum}   \stopitem
  \startitem \type{\overbars{Lorem ipsum}      }\hfill\overbars{Lorem ipsum}      \stopitem
  \startitem \type{\underbars{Lorem ipsum}     }\hfill\underbars{Lorem ipsum}     \stopitem
  \startitem \type{\overstrikes{Lorem ipsum}   }\hfill\overstrikes{Lorem ipsum}   \stopitem
  \startitem \type{\understrikes{Lorem ipsum}  }\hfill\understrikes{Lorem ipsum}  \stopitem
  \startitem \type{\hiddenbar{Lorem ipsum}     }\hfill\hiddenbar{Lorem ipsum}     \stopitem
  \startitem \type{\outline{Lorem ipsum}       }\hfill\outline{Lorem ipsum}       \stopitem
  \startitem \type{\outlined{Lorem ipsum}      }\hfill\outlined{Lorem ipsum}      \stopitem
  \startitem \type{\backgroundbar{Lorem ipsum} }\hfill\backgroundbar{Lorem ipsum} \stopitem
  \startitem \type{\mathbackground{Lorem ipsum}}\hfill\mathbackground{Lorem ipsum}\stopitem
  \startitem \type{\undergraphic{Lorem ipsum}  }\hfill\undergraphic{Lorem ipsum}  \stopitem
  \startitem \type{\underrandom{Lorem ipsum}   }\hfill\underrandom{Lorem ipsum}   \stopitem
  \startitem \type{\underrandoms{Lorem ipsum}  }\hfill\underrandoms{Lorem ipsum}  \stopitem
  \startitem \type{\underdash{Lorem ipsum}     }\hfill\underdash{Lorem ipsum}     \stopitem
  \startitem \type{\underdashes{Lorem ipsum}   }\hfill\underdashes{Lorem ipsum}   \stopitem
  \startitem \type{\underdot{Lorem ipsum}      }\hfill\underdot{Lorem ipsum}      \stopitem
  \startitem \type{\underdots{Lorem ipsum}     }\hfill\underdots{Lorem ipsum}     \stopitem
  \startitem \type{\underround{Lorem ipsum}    }\hfill\underround{Lorem ipsum}    \stopitem
  \startitem \type{\bigunderdot{Lorem ipsum}   }\hfill\bigunderdot{Lorem ipsum}   \stopitem
\stoppoints
  
\Topic{列表、表格、图片}

\Subject{列表环境示例}

\startbuffer
  \startpoints
    \item 这是无序列表
      \startpoints
        \item 这是无序列表
        \startpoints
          \item 这是无序列表
        \stoppoints
      \stoppoints
  \stoppoints
  
  \doloopoverlist
   {shadowedsquare,shadowedcircle,shadowedsquares,shadowedcircles,
    bullet,mp:bullet,
    circle,mp:circle,
    diamond,blackdiamond,asterisk,
    square,mp:square,blacksquare,mp:squares}
   {\startpoints[\recursestring]\item 这是无序列表\stoppoints}  
\stopbuffer

\typebuffer\blank

\startcolumns[n=2]
\getbuffer
\stopcolumns

\page
\null

\startbuffer
  \startordinals
    \item 这是有序列表
      \startordinals
        \item 这是有序列表
          \startordinals
            \item 这是有序列表
          \stopordinals
     \stopordinals
  \stopordinals
  
  \doloopoverlist
   {Co:num,CO:num,Cr:num,CR:num,Cs:num,CS:num,
   Co:Alph,CO:Alph,Cr:Alph,CR:Alph,Cs:Alph,CS:Alph,
   Co:alph,CO:alph,Cr:alph,CR:alph,Cs:alph,CS:alph,
   Co:cn,CO:cn,Cr:cn,CR:cn,Cs:cn,CS:cn,
   Co:hira,CO:hira,Cr:hira,CR:hira,Cs:hira,CS:hira,
   Co:kata,CO:kata,Cr:kata,CR:kata,Cs:kata,CS:kata,}
   {\startordinals[\recursestring,nostopper]\item 这是有序列表\stopordinals} 
\stopbuffer

\typebuffer\blank

\startcolumns[n=3]
\getbuffer
\stopcolumns

\Subject{表格环境}

\movesidefloat[y=-.3\baselineskip]
\startplacefigure [location={none,left,high,low},title=figure]
  \externalfigure [][height=6\baselineskip]
\stopplacefigure

\startplacefullfig [title=fullwidth figure]
  \externalfigure [][width=\dimexpr\textwidth+\rightmargintotal\relax,height=3\baselineskip]
\stopplacefullfig

\startplacetextfig [title={text figure},location={high,low}]
  \externalfigure [][height=4\baselineskip,width=1tw]
\stopplacetextfig

\startplacemarginfig [title=margin figure,location={outermargin}]
  \externalfigure [][width=\rightmarginwidth,height=5\baselineskip]
\stopplacemarginfig

\startplacefullfig [location=none]
   \startsubfloatnumbering
     \startfloatcombination [width=\dimexpr\textwidth+\rightmargintotal\relax,nx=3,distance=1mm,after=,inbetween=]
       \startplacefigure [title=Left]   \externalfigure \stopplacefigure
       \startplacefigure [title=Middle] \externalfigure \stopplacefigure
       \startplacefigure [title=Right]  \externalfigure \stopplacefigure
     \stopfloatcombination
   \stopsubfloatnumbering
\stopplacefullfig

\startbuffer
  \startlocalfootnotes
    \startxtable
        \startxrow[header]
          \startxcell 我走在命运为我规定的路上， 虽然我并不愿意走在这条路上， 
                      但是我除了满腔悲愤的走在这条路上别无选择\footnote{First} \stopxcell
          \startxcell 你今天是一个孤独的怪人，你离群索居，总有一天你会成为一个民族！\footnote{Second} \stopxcell
        \stopxrow
      \startxrow[body]
        \startxcell 我走在命运为我规定的路上， 虽然我并不愿意走在这条路上， 
                    但是我除了满腔悲愤的走在这条路上别无选择\footnote{First} \stopxcell
        \startxcell 你今天是一个孤独的怪人，你离群索居，总有一天你会成为一个民族！\footnote{Second} \stopxcell
      \stopxrow
      \startxrow[footer]
        \startxcell 受苦的人，没有悲观的权利。一个受苦的人，如果悲观了，就没有了面对现实的勇气，
                    也没有了与苦难抗争的力量，结果是他将受到更大的苦。\footnote{Third} \stopxcell
        \startxcell 所谓高贵的灵魂，即对自己怀有敬畏之心。\footnote{Fourth}  \stopxcell
      \stopxrow
      \startxrow[lastrow]
        \startxcell 受苦的人，没有悲观的权利。一个受苦的人，如果悲观了，就没有了面对现实的勇气，
                    也没有了与苦难抗争的力量，结果是他将受到更大的苦。\footnote{Third} \stopxcell
        \startxcell 所谓高贵的灵魂，即对自己怀有敬畏之心。\footnote{Fourth}  \stopxcell
      \stopxrow
    \stopxtable
    {\setupnotation[footnote][alternative={text},distance=.2em,after=\hskip1.5em]%
     \hbox{\placelocalfootnotes}}
  \stoplocalfootnotes
\stopbuffer

\movesidefloat[y=.3\baselineskip]
\startplacetable [location={none,left,high,low}]
  \getbuffer
\stopplacetable

\startplacefulltab [title=fullwidth table]
  \getbuffer
\stopplacefulltab

\startplacetexttab [title={text table},location={high,low,middle}]
    \getbuffer
\stopplacetexttab

\startplacefulltab [location=none]
   \startsubfloatnumbering
     \startfloatcombination [width=\dimexpr\textwidth+\rightmargintotal\relax,nx=3,distance=1mm,after=,inbetween=]
       \startplacetable [title=Left]   \getbuffer \stopplacetable
       \startplacetable [title=Middle] \getbuffer \stopplacetable
       \startplacetable [title=Right]  \getbuffer \stopplacetable
     \stopfloatcombination
   \stopsubfloatnumbering
\stopplacefulltab

\startplacetexttab[title=name of table]
  \getbuffer
\stopplacetexttab

\Topic{双栏标题框组件测试}

\Subject{基础功能测试}

\subsection{完整双栏标题框（默认）}

\startDualHeaderText[number={No.001},title={这是标题}]
这是一个完整的双栏标题框示例。左侧显示编号，右侧显示标题。
\stopDualHeaderText

\subsection{仅显示编号}

\startDualHeaderText[number={A-123},showtitle=0]
这个框只显示编号，不显示标题。
\stopDualHeaderText

\subsection{仅显示标题}

\startDualHeaderText[title={仅标题示例},shownumber=0]
这个框只显示标题，不显示编号。
\stopDualHeaderText

\subsection{无编号和标题}

\startDualHeaderText[number=,title=,shownumber=0,showtitle=0,toffset=0.5\bodyfontsize]
当编号和标题都为空时，不显示标题栏。需要手动调整 toffset 参数。
\stopDualHeaderText

\Subject{背景样式测试}

\subsection{simple 样式（默认）}

\startDualHeaderText[
  number=,
  title=,
  shownumber=0,
  showtitle=0,
  toffset=0.5\bodyfontsize,
  bodybackground=simple
]
这是默认的简单样式，带完整边框。
\stopDualHeaderText

\startDualHeaderText[
  number=,
  title=,
  shownumber=0,
  showtitle=0,
  toffset=0.5\bodyfontsize,
  bodybackground=simple,
  bodybgcolor=softthemecolor,
]
这是默认的简单样式，带完整边框。
\stopDualHeaderText

\subsection{symbolbox 样式}

\startDualHeaderText[
  number={Symbol},
  title={符号框样式},
  bodybackground=symbolbox,
  symbol=\usersymbolparameter{Frame}
]
symbolbox 样式会在右下角显示一个符号。
\stopDualHeaderText

\subsection{charbox 样式}

\startDualHeaderText[
  number={Char},
  title={字符图案样式},
  bodybackground=charbox,
  character={$\star$}
]
charbox 样式会在背景显示旋转45度的字符图案。
\stopDualHeaderText

\subsection{framebox 样式}

\startDualHeaderText[
  number={Char},
  title={字符图案样式},
  bodybackground=framebox,
]
framebox 样式会在背景显示棋盘格图案。
\stopDualHeaderText

\subsection{shadowbox 样式（自定义阴影色）}

\startDualHeaderText[
  number={Shadow},
  title={阴影框样式},
  bodybackground=shadowbox,
  shadowcolor=lightthemecolor,
]
shadowbox 样式会在右下角显示阴影效果。这里使用了自定义的阴影颜色。
\stopDualHeaderText

\subsection{shadowbox 样式（自动阴影色）}

\startDualHeaderText[
  number={Auto},
  title={自动阴影色},
  bodybackground=shadowbox,
  shadowcolor=white,
]
当 shadowcolor 为 white 时，自动使用边框颜色和白色的混合色。
\stopDualHeaderText

\subsection{barbox 样式}

\startDualHeaderText[
  number={Bar},
  title={条形样式},
  bodybackground=barbox,
  leftframe=1,
]
barbox 样式会在左侧显示一个色条。
\stopDualHeaderText

\Subject{边框控制测试}

\subsection{仅左边框}

\startDualHeaderText[
  shownumber=0,
  showtitle=0,
  leftframe=1,
  rightframe=0,
  topframe=0,
  bottomframe=0,
  toffset=0.5\bodyfontsize,
]
只显示左边框。
\stopDualHeaderText

\subsection{仅左右边框}

\startDualHeaderText[
  shownumber=0,
  showtitle=0,
  leftframe=1,
  rightframe=1,
  topframe=0,
  bottomframe=0,
  toffset=0.5\bodyfontsize,
]
只显示左右边框。
\stopDualHeaderText

\subsection{无边框}

\startDualHeaderText[
  shownumber=0,
  showtitle=0,
  leftframe=0,
  rightframe=0,
  topframe=0,
  bottomframe=0,
  toffset=0.5\bodyfontsize,
  bodybgcolor=softthemecolor,
]
不显示任何边框。
\stopDualHeaderText

\Subject{宽度比例测试}

\subsection{1:2 比例}

\startDualHeaderText[
  number={No.},
  title={标题宽度是编号的两倍},
  numberratio=1,
  titleratio=2
]
标题的宽度是编号的两倍。
\stopDualHeaderText

\subsection{1:3 比例}

\startDualHeaderText[
  number={A},
  title={标题宽度是编号的三倍},
  numberratio=1,
  titleratio=3
]
标题的宽度是编号的三倍。
\stopDualHeaderText

\Subject{圆角测试}

\subsection{圆角效果}

\startDualHeaderText[
  number={圆角},
  title={圆角效果演示},
  radius=10pt
]
使用 radius 参数设置圆角半径。标题栏背景色会自动遮盖圆角区域的空白。
\stopDualHeaderText

\subsection{大圆角}

\startDualHeaderText[
  number={大圆角},
  title={大圆角效果},
  radius=15pt
]
更大的圆角半径。
\stopDualHeaderText

\subsection{无圆角}

\startDualHeaderText[
  number={无圆角},
  title={无圆角效果},
  radius=0pt
]
默认无圆角。
\stopDualHeaderText

\Subject{标题栏背景色测试}

\subsection{自定义标题栏背景色}

\startDualHeaderText[
  number={自定义},
  title={自定义标题栏背景色},
  titlebgcolor=white,
  numberbgcolor=blue,
]
可以自定义标题栏的背景色。
\stopDualHeaderText

\subsection{透明标题栏背景}

\startDualHeaderText[
  number={透明},
  title={透明标题栏背景},
  titlebgcolor=transparentcolor,
  bodybgcolor=transparentcolor,
]
标题栏背景可以设置为透明。
\stopDualHeaderText

\Subject{颜色测试}

\subsection{自定义颜色}

\startDualHeaderText[
  number={自定义},
  title={自定义颜色示例},
  numberbgcolor=red,
  titlebgcolor=blue,
  bodybgcolor=yellow,
  framecolor=red,
  numbercolor=white,
  titlecolor=white,
]
可以自定义各种颜色。
\stopDualHeaderText

\Subject{组合效果测试}

\subsection{阴影 + 圆角}

\startDualHeaderText[
  number={组合},
  title={阴影+圆角},
  bodybackground=shadowbox,
  radius=8pt,
  shadowcolor=lightthemecolor,
]
阴影和圆角的组合效果。
\stopDualHeaderText

\subsection{左侧条 + 无标题}

\startDualHeaderText[
  number={组合},
  bodybackground=barbox,
  leftframe=1,
  showtitle=0,
]
左侧色条 + 无标题的组合效果。
\stopDualHeaderText

\subsection{字符图案 + 圆角}

\startDualHeaderText[
  number={组合},
  title={字符图案+圆角},
  bodybackground=charbox,
  character={$\heartsuit$},
  radius=10pt,
]
字符图案和圆角的组合效果。
\stopDualHeaderText

\subsection{仅左边框模拟 blockbox}

\startDualHeaderText[
  number={模拟},
  title={模拟 blockbox},
  leftframe=1,
  rightframe=0,
  topframe=0,
  bottomframe=0,
  rulethickness=2pt
]
通过设置仅显示左边框，可以模拟 blockbox 的效果。
\stopDualHeaderText

\Subject{titlepos=0（顶部标题）}

\subsection{完整双栏标题框}

\startDualHeaderText[
  number={No.1},
  title={顶部标题},
  titlepos=0,
]
这是 titlepos=0 的效果，标题显示在顶部，占用正文空间。
\stopDualHeaderText

\subsection{仅标题}

\startDualHeaderText[
  shownumber=0,
  title={仅标题},
  titlepos=0,
]
只显示标题，不显示编号。
\stopDualHeaderText

\subsection{仅编号}

\startDualHeaderText[
  shownumber=1,
  showtitle=0,
  number={No.2},
  titlepos=0,
]
只显示编号，不显示标题。
\stopDualHeaderText

\subsection{标题和编号（左对齐）}

\startDualHeaderText[
  number={No.3},
  title={左对齐标题},
  titlepos=0,
  titlealign=1,
]
标题左对齐，编号右对齐。
\stopDualHeaderText

\subsection{标题和编号（右对齐）}

\startDualHeaderText[
  number={No.4},
  title={右对齐标题},
  titlepos=0,
  titlealign=3,
]
标题右对齐，编号左对齐。
\stopDualHeaderText

\Subject{组合效果测试}

\subsection{titlepos=0 + 圆角}

\startDualHeaderText[
  number={No.5},
  title={圆角顶部标题},
  titlepos=0,
  radius=5pt,
]
titlepos=0 与圆角效果的组合。
\stopDualHeaderText

\subsection{titlepos=0 + 阴影}

\startDualHeaderText[
  number={No.6},
  title={阴影顶部标题},
  titlepos=0,
  bodybackground=shadowbox,
]
titlepos=0 与阴影效果的组合。
\stopDualHeaderText

\subsection{titlepos=0 + 自定义颜色}

\startDualHeaderText[
  number={No.7},
  title={自定义颜色},
  titlepos=0,
  numberbgcolor=blue,
  titlecolor=red,
  numbercolor=white,
]
titlepos=0 与自定义颜色的组合。
\stopDualHeaderText

\Subject{其他环境}

\startdescription{人性，太人性的}
人类的生命，不能以时间长短来衡量，心中充满爱时，刹那即为永恒！ ——尼采
完全不谈自己是一种甚为高贵的虚伪。 
\stopdescription

\starthangdescr{Description}
\input knuthmath
\stophangdescr

\starttopdescr{Description}
\input knuthmath
\stoptopdescr

\starttips{Description}
\input knuthmath
\stoptips

\startexcursus{Description}
\input knuthmath
\stopexcursus

\startremark
\input knuthmath
\stopremark

\Subject{和教学相关的一些环境}

\question{Lorem ipsum}

\question{Lorem ipsum\choice{Lorem ipsum,Lorem ipsum,Lorem ipsum,Lorem ipsum}}

\question{Lorem ipsum\problem{\pitem{Lorem ipsum}\pitem{Lorem ipsum}\pitem{Lorem ipsum}}}

\startlexicalentry[class=名,accent=1]{太陽}
  \startlexicallists
    \item 太陽系の中心をなし、地球にもっとも近い恒星。お日さま。\antonym{太陰}。
    \item 光明・希望の象徴。「心の\alias」
      \startlexicallists
        \item 太陽系の中心をなし、地球にもっとも近い恒星。お日さま。\example{太陰}。
        \item 光明・希望の象徴。「心の\alias」
          \startlexicallists
            \item 太陽系の中心をなし、地球にもっとも近い恒星。お日さま。\refer{太陰}。
            \item 光明・希望の象徴。「心の\alias」
              \startlexicallists
                \item 太陽系の中心をなし、地球にもっとも近い恒星。お日さま。\junction{太陰}。
                \item 光明・希望の象徴。「心の\alias」
              \stoplexicallists
          \stoplexicallists
      \stoplexicallists
  \stoplexicallists
\stoplexicalentry

\startlexicalsyntax[class=名,accent=1]{太陽}
  \lexicalgrammar 太陽系の中心をなし、地球にもっとも近い恒星。\par
  \lexicalgrammar 太陽系の中心をなし、地球にもっとも近い恒星。\par
  \startlexicalgrammar
    太陽系の中心をなし、地球にもっとも近い恒星。
    \startlexicallists
      \item 太陽系の中心をなし、地球にもっとも近い恒星。お日さま。\lexicalemph{太陰}。
      \item 光明・希望の象徴。「心の\alias」
    \stoplexicallists
  \stoplexicalgrammar

  \startlexicalexample ワインをは飲みますが、ビールをはのみません。\stoplexicalexample
  \startlexicalexample ワインをは飲みますが、ビールをはのみません。\stoplexicalexample
  \startlexicalexample
    ワインをは飲みますが、ビールをはのみません。
    \startlexicallists
      \item 私はワインは飲みますが、ビールは飲みません。
      \item 映画は見ませんが、ドラマは見ます。
      \item 豚肉は食べませんが、牛肉や鶏肉はよく食べます。
    \stoplexicallists
  \stoplexicalexample
  
  \startlexicalmeaning{Aの状態ではなくなるまでにBを完了する}\\
    \junction{動詞辞書形・テイル形・ナイ形}\\
    \junction{形容詞・形容動詞の連体形（現在肯定形・現在否定形）}\\
    \junction{名詞＋の}
    \startlexicallists
      \item 明るいうちに、帰ったほうがいい。
      \item 元気なうちに、世界一周旅行がしたい。
    \stoplexicallists
  \stoplexicalmeaning

  \startlexicaljunction
    太陽系の中心をなし、地球にもっとも近い恒星。
    \startlexicallists
      \item 明るいうちに、帰ったほうがいい。
      \item 元気なうちに、世界一周旅行がしたい。
    \stoplexicallists
  \stoplexicaljunction
  
  \startlexicalcomparing
    太陽系の中心をなし、地球にもっとも近い恒星。
    \startlexicallists
      \item 明るいうちに、帰ったほうがいい。
      \item 元気なうちに、世界一周旅行がしたい。
    \stoplexicallists
  \stoplexicalcomparing
  
  \startlexicalsummary
    太陽系の中心をなし、地球にもっとも近い恒星。
    \startlexicallists
      \item 明るいうちに、帰ったほうがいい。
      \item 元気なうちに、世界一周旅行がしたい。
    \stoplexicallists
  \stoplexicalsummary
\stoplexicalsyntax

\stoptext
